\documentclass[sigconf,anonymous,review]{acmart}

\usepackage{amsmath}
\usepackage{booktabs}

\setcopyright{none}
\acmDOI{}
\acmISBN{}
\newcommand{\system}{EPM}

\title{Supplementary Material for ``Proactive Memory for Event-Driven LLM Agents''}

\begin{document}
\maketitle

\section{Supplementary Entity-Composition Sensitivity Check}

We use a cluster bootstrap solely as a descriptive check of entity composition,
constrained to preserve every aggregate confusion matrix.  The frozen TP and FP
totals for each benchmark and method are allocated to the
constituent entities in proportion to their positive and negative reference
support, using largest-remainder allocation so that every aggregate confusion
matrix is preserved exactly.  We then resample the nine HeartSteps participants
and four CASAS homes with replacement within benchmark, recompute micro F1, and
repeat this procedure 100,000 times (fixed seed 20260824).

\begin{table}[h]
\centering
\caption{Micro F1 and 2.5--97.5 percentile ranges under the allocation model.}
\label{tab:supp_method_ci}
\small
\begin{tabular}{lrr}
\toprule
Method & Micro F1 & 2.5--97.5 range \\
\midrule
Source-only & 38.8 & [34.6, 43.0] \\
Source+Working & 40.7 & [36.7, 44.7] \\
Full \system{} & 51.7 & [47.1, 56.1] \\
PA-Prompt & 41.2 & [35.6, 46.7] \\
\bottomrule
\end{tabular}
\end{table}

\begin{table}[h]
\centering
\caption{Percentile ranges for the micro-F1 advantage of Full \system{} under the allocation model.}
\label{tab:supp_delta_ci}
\small
\begin{tabular}{lrr}
\toprule
Comparison & $\Delta$F1 & 2.5--97.5 range \\
\midrule
Full \system{} vs. PA-Prompt & 10.5 & [8.4, 12.5] \\
Full \system{} vs. Source-only & 12.9 & [12.1, 13.4] \\
Full \system{} vs. Source+Working & 11.0 & [10.3, 11.5] \\
\bottomrule
\end{tabular}
\end{table}

All three percentile ranges remain above zero under changes in the sampled
entity composition.  The ranges are conditional on the stated proportional
allocation model.  This check is not a hypothesis test, does not yield
$p$-values, and does not estimate within-entity model-sampling variance.

\section{Deterministic Evidence Validation}

The deterministic layer verifies whether memory evidence is admissible.  It
does not decide whether a proactive service should be sent and never reads a
reference label.

\subsection{Daily fact validation}

For a proposed daily episode $d$, the validator recursively extracts scalar
leaves only from its structured \texttt{observed},
\texttt{context-before}, and \texttt{outcome} fields.  Free-text summaries and
inferred fields are excluded.  Scalar values are converted to a stable typed
representation.  A claimed leaf is grounded when its typed value equals a leaf
in a cited source event and its field path is either identical to, or ends with,
the claimed path.  Every citation must refer to the same entity and completed
local day and must fall within the episode interval.  Unknown citations,
out-of-interval citations, and unsupported structured leaves cause rejection.
Only episodes satisfying all checks enter Working Memory.

\subsection{Weekly historical replay}

A long-term proposal has a trigger context $g$, an expected pattern $p$, an
optional weekday restriction, and an expected local-time window.  Replay uses
only the validated structured daily facts.  For each applicable day, missing or
unobservable data are censored.  A non-censored day becomes eligible when $g$
is observed.  It is fulfilled when $p$ occurs after the first trigger and inside
the proposed time window; otherwise it is violated.  We compute
\begin{align}
\mathrm{confidence} &=
\frac{|\mathrm{fulfilled}|}{|\mathrm{fulfilled}|+|\mathrm{violated}|},\\
\mathrm{lift} &=
\frac{\mathrm{confidence}}{\mathrm{base\ rate}},
\end{align}
where the base rate is the frequency of $p$ in the same window over observable
days without conditioning on the trigger.  A proposal becomes Active when it
has at least three supporting days, confidence at least $0.5$, and lift at least
$1.0$.  A well-formed, source-linked proposal below these thresholds remains
Candidate and is replayed again after later completed weeks.  Malformed
conditions, unknown evidence identifiers, and conditions that cannot be
reconstructed online fail closed.  Candidate proposals are not exposed to the
Decision LLM.

\section{Online Relevance and Episode Control}

For a current event, only Active long-term trajectories whose trigger, weekday,
and expected window match are considered.  The default
lead and grace intervals are both 30 minutes.  Each match receives one of five
relations: \texttt{not-started}, \texttt{pending}, \texttt{fulfilled},
\texttt{deviated}, or \texttt{censored}.  Fulfilled and censored items are not
placed in the Evidence Card.  Remaining items are ranked by relation
(deviated, pending, not-started), status, replay confidence, and temporal
distance; at most three are shown to the Decision LLM.

The Episode Controller acts only after the Decision LLM proposes a Push.
For HeartSteps, the episode signature contains availability, connectivity,
snooze and transit state, recognized activity, and location category.  For
CASAS, it contains the activity family and observed room set.  Consecutive
events with the same signature on the same local day belong to one episode.
A day boundary or signature change starts a new episode.  CASAS additionally
distinguishes onset, confirmed, and sustained evidence phases, allowing a
condition with materially stronger support to be reconsidered.

The intervention identity combines the episode, service target or intent, and
the long-term evidence identifiers used by the decision.  If the same identity
has already produced a Push, a repeated candidate is converted to Abstain.
A changed episode signature, a new day, a stronger evidence phase, a changed
service, or changed trajectory evidence creates a new identity and permits a
fresh Decision-LLM judgment.  The controller cannot turn Abstain into Push.

The causal execution order is therefore: read strictly prior memory; stage the
current event; build the Evidence Card and episode identity; ask the Decision
LLM; apply duplicate suppression only to a candidate Push; commit the event;
and finally run any daily or weekly maintenance that has become due.

\subsection{Operational reproducibility checklist}

Table~\ref{tab:operational_contracts} summarizes the deterministic contracts
used by the frozen evaluation.  These contracts govern evidence admission and
duplicate suppression; none of them directly predicts Push or reads a reference
label.

\begin{table}[h]
\centering
\caption{Frozen operational contracts.}
\label{tab:operational_contracts}
\small
\begin{tabular}{p{0.29\columnwidth}p{0.62\columnwidth}}
\toprule
Question & Frozen contract \\
\midrule
Daily fact parsing & Typed scalar leaves from structured observed, prior-context, and outcome fields only. \\
Fact agreement & Same typed value and identical or suffix-matching field path in a cited same-entity event. \\
Active threshold & At least three support days, confidence $\geq0.5$, and lift $\geq1.0$. \\
Candidate lifecycle & Replay well-formed insufficient proposals; reject malformed or unreconstructable proposals. \\
Long-term relevance & Match trigger, weekday, and time window; then rank unresolved relations by evidence strength and distance. \\
Evidence Card cap & At most three relevant long-term trajectories. \\
Episode identity & Domain state signature plus local-day boundary; CASAS also tracks evidence phase. \\
Close or reopen & Suppress a repeated intervention identity; reopen on a new day, changed signature, stronger evidence, changed service, or changed trajectory evidence. \\
\bottomrule
\end{tabular}
\end{table}

\subsection{Prompt and schema contract}

Source-only, Source+Working, and Full \system{} share one decision instruction:
return Push only when the current event or visible memory evidence supports a
concrete, timely, feasible, and low-risk service opportunity; otherwise return
Abstain.  The instruction prohibits unsupported health, danger, travel, and
preference claims.  The decision schema contains the decision, service intent,
optional message, and cited evidence identifiers.  The Episode Controller reads
this output only after a candidate Push has been proposed.

Daily construction returns an episode identifier, local interval, structured
\texttt{observed}, \texttt{context-before}, and \texttt{outcome} fields, and
source-event citations.  Long-term construction returns a trajectory
identifier, trigger, expected pattern, optional weekday, local-time window, and
daily-evidence citations.  These schemas delimit what the deterministic
validators may accept; free-text explanations are never replayed as facts.

\end{document}
